% Drawing 79/08-1 -- folder 79, batch 08, archive page 155 (PDF page 156).
% Transcription: en bas à droite, dans l'encadré biffé : un point $s$ relié à deux points $t_1$ et $t_0$ d'une droite $D$
% Kind: point joined to two points of a line. Method 2 (schematic redraw from the facsimile). Status: ok.
% Notes: Inside the struck box (its frame lines are not drawn): a point s joined to t_1 and t_0, both on the line D. Strokes not struck through individually.
% Coordinates: 1 unit = 1 mm = 1/100 of the crop width (bbox in page fractions: [0.67, 0.78, 0.85, 0.87]).
% Redrawn by Claude (claude-opus-5-5), issue #73, 2026-11-07; not checked by a human; not compiled (no TeX in the session).
% Rights: reproduces part of the work -- behind RENDER_TRANSCRIPTS, never published by itself.
\documentclass[tikz,border=3pt]{standalone}
\usepackage{amsmath,amssymb}
\usetikzlibrary{patterns,arrows.meta}
\tikzset{pencil/.style={black!45}, struck/.style={black!35}, mark/.style={circle,draw,inner sep=1.3pt}, pt/.style={circle,fill,inner sep=0.9pt}, opt/.style={circle,draw,fill=white,inner sep=1.1pt}}
\begin{document}
\begin{tikzpicture}[x=1mm,y=1mm,line cap=round,line join=round,font=\scriptsize]
\draw (90,-5) -- (57,-58);
\draw (25,-31) -- (83,-12);
\draw (25,-31) -- (65,-44);
\node[pt] at (25,-31) {};
\node[pt] at (83,-12) {};
\node[pt] at (65,-44) {};
\node[inner sep=1pt] at (20,-29) {\scriptsize $s$};
\node[inner sep=1pt] at (79,-8) {\scriptsize $t_1$};
\node[inner sep=1pt] at (69,-47) {\scriptsize $t_0$};
\node[inner sep=1pt] at (93,-17) {\scriptsize $D$};
\end{tikzpicture}
\end{document}
